% =====================================================================
%  BAB III METODOLOGI
% =====================================================================
\chapter{METODOLOGI}
\label{bab:metodologi}

\petunjuk{Cara penelitian dilakukan, cukup rinci agar dapat direplikasi
(Pedoman TA 7.3). \tahap{\WAJIB}{\WAJIB{} (final)}}

\section{Alur Penelitian}
\petunjuk{Diagram alir tahapan (studi literatur, pengumpulan data,
perancangan, implementasi, pengujian, analisis) dan penjelasan naratif tiap
tahap.}

\begin{figure}[H]
\centering
% Ganti kotak berikut dengan: \includegraphics[width=0.8\linewidth]{gambar/alur-penelitian.png}
\fbox{\parbox[c][5cm][c]{0.7\linewidth}{\centering
  \isian{Sisipkan diagram alir penelitian}}}
\caption{Diagram Alir Penelitian}
\label{fig:alur-penelitian}
\end{figure}

\isian{Narasi tiap tahap pada Gambar~\ref{fig:alur-penelitian}.}

\section{Pengumpulan Data}
\petunjuk{Sumber data, karakteristik, metode pengumpulan, praproses, dan
pembagian data latih/validasi/uji tanpa kebocoran data. Sertakan izin, etika,
dan anonimisasi bila memakai data pribadi (UU No.~27 Tahun 2022). Uraikan
secara naratif: sumber dan lisensi/izin, jumlah sampel dan atribut, metode
pengumpulan, praproses, dan pembagian data.}

\isian{Tulis uraian dalam bentuk paragraf.}

\section{Perancangan Sistem/Model}
\petunjuk{Arsitektur, pemilihan algoritma beserta justifikasi, hiperparameter,
dan spesifikasi teknis. Untuk perangkat lunak sertakan UML (\textit{use case},
\textit{class}, \textit{sequence}) atau DFD.}

\isian{Tulis uraian dalam bentuk paragraf.}

\section{Lingkungan Eksperimen}
\petunjuk{Wajib untuk reprodusibilitas. Uraikan secara naratif: CPU / GPU /
RAM, sistem operasi, bahasa pemrograman dan versi, \textit{framework/library}
dan versi, platform \textit{cloud} (bila ada), \textit{seed} acak, dan
repositori kode.}

\isian{Tulis uraian dalam bentuk paragraf.}

\section{Metode Pengujian dan Evaluasi}
\petunjuk{Metrik evaluasi beserta rumus (nomor persamaan rata kanan) dan
skenario uji, misalnya \textit{k-fold cross-validation}, uji pada dataset
terpisah, \textit{stress test}, atau uji usabilitas.}

\isian{Tulis uraian dalam bentuk paragraf.} Contoh penulisan persamaan
bernomor seperti Persamaan~\eqref{eq:akurasi}:
\begin{equation}
  \text{Akurasi} = \frac{TP + TN}{TP + TN + FP + FN}
  \label{eq:akurasi}
\end{equation}
dengan $TP$, $TN$, $FP$, dan $FN$ berturut-turut menyatakan
\textit{true positive}, \textit{true negative}, \textit{false positive}, dan
\textit{false negative}.
